% ============================================================
% PREAMBLE.TEX — ELN / bookdown + xelatex
% Versão robusta: rodapé padronizado, sem duplicações,
% com suporte a química (mhchem), unidades (siunitx), CJK fallback.
% ============================================================

% ─────────────────────────────────────────────────────────────
% 1. FONTES E IDIOMA
% ─────────────────────────────────────────────────────────────
\usepackage{fontspec}

% Fonte monoespaçada para código
\setmonofont{DejaVu Sans Mono}[Scale=0.85]

% Fonte principal: sans-serif legível (DejaVu Sans cobre acentos PT + símbolos comuns)
\setmainfont{DejaVu Sans}[Scale=0.92]
\setsansfont{DejaVu Sans}[Scale=0.92]

% Fallback CJK (chinês) — requer fonts-noto-cjk instalado no sistema
% \usepackage{xeCJK}
% \setCJKmainfont{Noto Serif CJK SC}
% \setCJKsansfont{Noto Sans CJK SC}
% \setCJKmonofont{Noto Sans Mono CJK SC}

\usepackage{microtype}
\usepackage[portuguese,main=english]{babel}

% Entrelinha 1,5
\usepackage{setspace}
\onehalfspacing

% ─────────────────────────────────────────────────────────────
% 2. GEOMETRIA E LAYOUT
% ─────────────────────────────────────────────────────────────
\usepackage{geometry}
\geometry{
  a4paper,
  margin=1.2cm,
  headheight=14pt,
  headsep=0.4cm,
  footskip=0.9cm,   % aumentado para caber o rodapé de 3 colunas
  twoside=false
}

\usepackage{pdflscape}
\usepackage{afterpage}

% ─────────────────────────────────────────────────────────────
% 3. TABELAS (kableExtra / booktabs / longtable)
% ─────────────────────────────────────────────────────────────
\usepackage{booktabs}
\usepackage{longtable}
\usepackage{array}
\usepackage{multirow}
\usepackage{makecell}
\usepackage{colortbl}
\usepackage{xcolor}
\usepackage{tabularx}
\usepackage{tabulary}
\usepackage{threeparttable}

\newcolumntype{Y}{>{\raggedright\arraybackslash}X}
\newcolumntype{Z}{>{\centering\arraybackslash}X}

% ─────────────────────────────────────────────────────────────
% 4. GRÁFICOS E FIGURAS
% ─────────────────────────────────────────────────────────────
\usepackage{graphicx}
\usepackage{grfext}
\usepackage{truncate}

\graphicspath{%
  {./figs/}%
  {../figs/}%
  {Imagens_Caderno_2026/}%
}

% ─────────────────────────────────────────────────────────────
% 5. HIPERLINKS
% ─────────────────────────────────────────────────────────────
\usepackage[
  colorlinks=true,
  linkcolor=blue!60!black,
  urlcolor=blue!60!black,
  citecolor=blue!60!black,
  filecolor=blue!60!black,
  linktoc=all,
  pdfborder={0 0 0},
  pdfpagemode=UseOutlines,
  bookmarks=true,
  bookmarksopen=true,
  bookmarksopenlevel=2
]{hyperref}

\usepackage{xurl}

% ─────────────────────────────────────────────────────────────
% 6. MATEMÁTICA, QUÍMICA E UNIDADES
% ─────────────────────────────────────────────────────────────
\usepackage{amsmath}
\usepackage{amssymb}
\usepackage{bm}
\usepackage{bbm}
\usepackage{dsfont}
\usepackage{mathrsfs}
\usepackage{eucal}
\usepackage{upgreek}
\usepackage{wasysym}
\usepackage{stmaryrd}

% Química: \ce{H2O}, \ce{H2SO4}, etc.
\usepackage[version=4]{mhchem}

% Unidades: \SI{5}{\micro\liter}, \SI{37}{\celsius}, \SI{12000}{rpm}
\usepackage{siunitx}
\sisetup{
  detect-all,
  group-separator={.},
  output-decimal-marker={,}
}

% ─────────────────────────────────────────────────────────────
% 7. LISTAS
% ─────────────────────────────────────────────────────────────
\usepackage{enumitem}
\setlist{nosep}
\setlist[itemize]{leftmargin=*}
\setlist[enumerate]{leftmargin=*}

% ─────────────────────────────────────────────────────────────
% 8. FORMATAÇÃO DE TEXTO
% ─────────────────────────────────────────────────────────────
\usepackage{textcomp}

\parindent=0mm
\parskip=1.5ex plus0.5ex minus0.5ex

% ─────────────────────────────────────────────────────────────
% 9. CÓDIGO / VERBATIM
% ─────────────────────────────────────────────────────────────
\usepackage{fvextra}

\fvset{
  breaklines=true,
  breakanywhere=true,
  frame=single,
  framesep=2mm,
  fontsize=\small
}

\usepackage{csquotes}

% ─────────────────────────────────────────────────────────────
% 10. FIGURAS E LEGENDAS
% ─────────────────────────────────────────────────────────────
\usepackage{float}
\usepackage{caption}
\usepackage{subcaption}

\floatplacement{figure}{H}
\floatplacement{table}{H}

\captionsetup{
  font=small,
  labelfont=bf,
  justification=raggedright,
  singlelinecheck=false
}

% ─────────────────────────────────────────────────────────────
% 11. PACOTES ÚTEIS
% ─────────────────────────────────────────────────────────────
\usepackage{adjustbox}
\usepackage{calc}
\usepackage{etoolbox}
\usepackage{lastpage}   % para \pageref{LastPage} no rodapé

% ─────────────────────────────────────────────────────────────
% 12. CORREÇÕES DE COMPATIBILIDADE
% ─────────────────────────────────────────────────────────────

% Evitar que tabelas/verbatim herdem entrelinha 1,5
\AtBeginEnvironment{longtable}{\singlespacing}
\AtBeginEnvironment{tabular}{\singlespacing}
\AtBeginEnvironment{tabularx}{\singlespacing}
\AtBeginEnvironment{verbatim}{\singlespacing}
\AtBeginEnvironment{Highlighting}{\singlespacing}

% ─────────────────────────────────────────────────────────────
% 13. APÊNDICES
% ─────────────────────────────────────────────────────────────
\usepackage[title]{appendix}
\usepackage{footnotehyper}
\usepackage{footnotebackref}

% ─────────────────────────────────────────────────────────────
% 14. RODAPÉ PADRONIZADO (único para todos os documentos)
% ─────────────────────────────────────────────────────────────
% Macros configuráveis. Cada Rmd pode sobrescrever no chunk de setup:
%   cat("\\renewcommand{\\footLeft}{POP-042 -- Rev.~3}")
%   cat("\\renewcommand{\\footCenter}{\\thepage\\ / \\pageref{LastPage}}")
%   cat("\\renewcommand{\\footRight}{\\today}")
% Se não sobrescrever, os fallbacks abaixo são usados.

\providecommand{\popCode}{Código}
\providecommand{\popVersion}{0}
\providecommand{\popGenerated}{--/--/----}
\providecommand{\popCommit}{unknown}

\providecommand{\footLeft}{\popCode\ -- Rev.~\popVersion}
\providecommand{\footCenter}{\thepage\ / \pageref{LastPage}}
\providecommand{\footRight}{Gerado em \popGenerated}

\usepackage{fancyhdr}
\pagestyle{fancy}
\fancyhf{}

% Cabeçalho: número de página à direita, secção à esquerda
\fancyhead[L]{\small\nouppercase{\leftmark}}
\fancyhead[R]{\small\thepage}

% Rodapé: três colunas com macros configuráveis
\fancyfoot[L]{\small\footLeft}
\fancyfoot[C]{\small\footCenter}
\fancyfoot[R]{\small\footRight}

\renewcommand{\headrulewidth}{0.4pt}
\renewcommand{\footrulewidth}{0.4pt}

% ─────────────────────────────────────────────────────────────
% 15. CONFIGURAÇÕES FINAIS
% ─────────────────────────────────────────────────────────────
\hyphenpenalty=10000
\exhyphenpenalty=10000
\Urlmuskip=0mu plus 1mu
\hfuzz=1pt
\vfuzz=1pt

% Forçar oneside (redundância defensiva)
\makeatletter
\@twosidefalse
\@mparswitchfalse
\makeatother
